% 13 · The benchmark — what we measured, and what it does not support
\section{THE BENCHMARK: WHAT WE MEASURED}
\label{sec:benchmark}

\deckslides{slides 46--47 for the benchmark and slide 48 for the lessons in short form.}

This chapter is the case study: the methods of
\S\ref{sec:kmeans}--\S\ref{sec:evoc} run on the data of \S\ref{sec:data},
scored by the protocol of \S\ref{sec:validation}. The order of business is the
benchmark grid, then the two tasks separately, then the levers, and finally the
list of things this benchmark cannot tell you.

\subsection{The design}

Three things vary, and nothing else:

\begin{itemize}
\item \textbf{Signal.} The 16 standardised abundances of \S\ref{sec:data};
  a 64- or 256-dimensional PCA projection of the APOGEE spectra
  (\S\ref{sec:spectral}); a 256-dimensional masked-autoencoder latent of the
  same spectra; the four kinematic observables
  (parallax, two proper motions, radial velocity); and concatenations of
  chemistry with kinematics.
\item \textbf{Method.} t-SNE, UMAP and EVoC, each followed by
  HDBSCAN$^*$ clustering where the method needs one, with the defaults of
  \textbf{Table~\ref{tab:config}}.
\item \textbf{Task.} \emph{Cluster-only separation}: cluster the known members
  among themselves and ask whether the clusters come out distinct, scored with
  \textbf{Equation~\ref{eq:merits}}. \emph{Field retrieval}: search the full
  sample for members and score with \textbf{Equation~\ref{eq:rp}}.
\end{itemize}

Three conventions keep the comparison honest, and all three are choices a
reader should check before believing any row:

\begin{enumerate}
\item \textbf{Same population.} Every arm in \textbf{Table~\ref{tab:honest}}
  is scored on the same 982 stars and the same 25 clusters. An earlier version
  of this table compared arms scored on 25 clusters against arms scored on 5,
  which is the mistake of \S\ref{sec:validation}.
\item \textbf{Seven seeds, reported as mean $\pm$ s.d.} The kinematics-only
  t-SNE row is deterministic, several other rows are not, and the difference
  is information rather than noise to be hidden.
\item \textbf{Two referees.} Field retrieval is scored against two independent
  definitions of ``member'': the kinematic selection of
  \S\ref{sec:data}, and the Simbad cluster memberships of the literature.
  Every conclusion below holds under both; the Simbad referee moves precision
  by $-0.01$ to $-0.08$ (the largest shifts are in the 64-dimensional PCA arms)
  and changes no ranking.
\end{enumerate}

\subsection{Task 1: separating the clusters from each other}

\textbf{Table~\ref{tab:honest}} is the central result of the workbook. Read it
by column first: the columns are methods, the rows are the signal the method
was given.

\begin{table}[htbp]
\centering
\small
\tabcolsep4pt
\begin{tabular}{@{}llccc@{}}
\toprule
signal & dim & t-SNE & UMAP & EVoC \\
\midrule
abundances                & 16  & $0.560 \pm 0.000$ & $0.578 \pm 0.011$ & $0.420 \pm 0.043$ \\
PCA of spectra            & 64  & $0.741 \pm 0.000$ & $0.754 \pm 0.007$ & $0.733 \pm 0.008$ \\
PCA of spectra            & 256 & $0.688 \pm 0.003$ & $0.685 \pm 0.014$ & $0.653 \pm 0.020$ \\
masked AE latent          & 256 & $0.740 \pm 0.002$ & $0.760 \pm 0.011$ & $0.714 \pm 0.036$ \\
kinematics only           & 4   & $0.942 \pm 0.000$ & $0.940 \pm 0.005$ & $0.876 \pm 0.012$ \\
\midrule
abundances, K-means ($K$ known) & 16 & \multicolumn{3}{c}{$0.449$ (single run)} \\
\bottomrule
\end{tabular}
\caption{Cluster-only homogeneity on the same 982 stars in 25 clusters, mean
$\pm$ standard deviation over seven seeds. Higher is better. The K-means row
is the classical baseline of \S\ref{sec:kmeans}, run on the 1\,002-row
member matrix with $K=25$ and $n_{\rm init}=10$ by
\texttt{article/scripts/kmeans\_baseline.py}: it is not an embedding, so it
has no column to sit in, and it is quoted for scale rather than for ranking.
The recovery fractions of the same arms, and the comparison with the published
result, are in \textbf{Table~\ref{tab:literature}}.}
\label{tab:honest}
\end{table}

Four things in that table are worth stating out loud.

\textbf{Chemistry is not nothing.} All three methods reach homogeneity well
above the trivial floor of $1/25 = 0.04$: the 16 abundances carry real cluster
information, and the embedding methods extract some of it. This is the
optimistic half of the chemical-tagging programme, reproduced on data that
post-date the paper that reported it.

\textbf{Chemistry is not enough.} The same methods on four kinematic
observables --- a quarter of the dimensions --- reach $0.94$. Nothing about
the abundance arms is competitive with kinematics, and no amount of embedding
repairs it: the four-dimensional signal is simply more informative about
which cluster a star belongs to.

\textbf{Dimension helps, up to a point.} Going from 16 abundances to the
64-dimensional PCA latent raises EVoC's homogeneity from $0.42$ to $0.73$ and
t-SNE's from $0.56$ to $0.74$. Going on to 256 principal components \emph{lower}
it --- $0.69$ for t-SNE, $0.65$ for EVoC --- which is the bias--variance trade
of \S\ref{sec:fundamentals} in its least forgiving form: the extra 192
directions carry variance, not signal, and every one of them is a direction in
which the field can become an apparent neighbour.

\textbf{One row of that table is not a measurement.} The abundances-only
t-SNE row is unstable under row order, as the audit of \S\ref{sec:validation}
showed: the same matrix gives $0.218$--$0.560$ depending on how the rows are
ordered, and the single-run baseline of \textbf{Table~\ref{tab:baseline}}
gives $0.292$ for what is nominally the same experiment. Any claim that rests
on the difference between $0.29$ and $0.56$ for t-SNE on raw abundances is
built on the input ordering.

\begin{table}[htbp]
\centering
\small
\tabcolsep5pt
\begin{tabular}{@{}lccc@{}}
\toprule
features & t-SNE & UMAP & EVoC \\
\midrule
abundances (16-d)          & 0.292 & 0.547 & 0.409 \\
abundances + kinematics    & \textbf{0.767} & 0.748 & 0.680 \\
\bottomrule
\end{tabular}
\caption{The single-run baseline of the DR19 re-run
(\texttt{cluster baseline}, no seed averaging, the full 1\,002-row member
matrix). It is included because it is the number our own documentation quotes,
and because comparing it with \textbf{Table~\ref{tab:honest}} is instructive:
concatenating kinematics lifts every method by $0.25$--$0.48$, and the
abundances-only t-SNE row differs between the two tables by more than the seed
spread of any other row.}
\label{tab:baseline}
\end{table}

\begin{table}[htbp]
\centering
\small
\tabcolsep5pt
\begin{tabular}{@{}llccc@{}}
\toprule
signal (excluding M~3) & t-SNE & UMAP & EVoC \\
\midrule
abundances (884 stars)              & $0.199$ (degenerate) & $0.581$ & $0.499$ \\
masked AE latent (884 stars)        & $0.711$ & $0.742$ & $0.658$ \\
\bottomrule
\end{tabular}
\caption{Ablation on the 982-star population: the globular cluster M~3
contributes $98$ of the stars ($884$ remain) --- or $154$ of the
1\,002 member rows of \textbf{Table~\ref{tab:clusters}} --- and it is the
largest globular in the sample, though not the largest cluster: M~67 supplies
$218$ stars and the Pleiades $113$. Remove M~3 and the abundances-only t-SNE arm does not
merely change, it collapses: two groups with $75\%$ of the stars in the
largest, zero stars labelled noise. The spectral latent is unaffected by the
same ablation. A result that depends on one cluster being present is not a
measurement of the method. The per-cluster counts of both populations are in
\texttt{results/population\_counts.csv}
(\texttt{python scripts/population\_counts.py --baseline-min-finite 9}).}
\label{tab:ablation}
\end{table}

\begin{figure}[htbp]\centering
  \fig[0.98]{benchmark_grid.png}
  \caption{The whole grid in one figure: method (rows) against signal
  (columns), for the two tasks. Read the left block for cluster-only
  separation and the right block for field retrieval, and notice that the
  two blocks do not agree on a winner --- no single row is best at both. This
  is the figure the lecture ends the benchmark on, and it is the reason
  \S\ref{sec:benchmark} refuses to name a best method.}
  \label{fig:benchgrid}
\end{figure}

\begin{figure}[htbp]\centering
  \fig[0.86]{headtohead_pca.png}
  \caption{The head-to-head on $64$-dimensional PCA spectra, cluster-only
  separation, seven seeds. The error bars are the seed spread: the masked-AE
  and PCA-64 arms overlap, which is the tie reported in
  \S\ref{sec:spectral} and the reason the spectral result is stated as
  ``spectra beat abundances, and do not beat PCA''.}
  \label{fig:headtohead}
\end{figure}

\subsection{Task 2: finding members in the field}

Now the operational question. Take 30\,107 stars (the 25 clusters plus a
stratified sample of the field), cluster them, and ask how many of the known
members were recovered and how much of what was returned was real.

\begin{table}[htbp]
\centering
\small
\tabcolsep4pt
\begin{tabular}{@{}llccc@{}}
\toprule
signal & method & recall & precision & kNN purity \\
\midrule
abundances (16-d) & t-SNE & 0.201 & 0.240 & 0.124 \\
abundances (16-d) & UMAP  & 1.000 & 0.001 & 0.075 \\
abundances (16-d) & EVoC  & 0.557 & 0.003 & --- \\
\midrule
PCA 64-d & t-SNE & 0.417 & 0.631 & 0.528 \\
PCA 64-d & UMAP  & 0.374 & 0.592 & 0.434 \\
PCA 64-d & EVoC  & 0.585 & 0.003 & --- \\
\midrule
masked AE 256-d & t-SNE & 0.418 & 0.619 & 0.504 \\
masked AE 256-d & UMAP  & 0.385 & 0.513 & 0.419 \\
masked AE 256-d & EVoC  & 0.255 & 0.005 & --- \\
\bottomrule
\end{tabular}
\caption{Field retrieval, all-sky, on the same 30\,107 stars and the same 25
clusters, scored against the kinematic membership (978 members). Recall and
precision are the best-overlap definition of \textbf{Equation~\ref{eq:rp}};
kNN purity is the parameter-free cohesion score of \S\ref{sec:knn}, and it is
undefined for EVoC because that arm returns no embedding to measure it in.
Scoring against the Simbad referee instead shifts every precision by
$-0.01$ to $-0.08$ and changes no ranking.}
\label{tab:field}
\end{table}

The three lessons here organise most of what the rest of this workbook says
about chemical tagging in practice.

\textbf{Recall without precision is a statement about the field, not about
the clusters.} UMAP on raw abundances recovers \emph{every} true member ---
at a precision of $0.001$. It has not found the clusters; it has labelled a
large part of the field as cluster, and the best-overlap protocol dutifully
reports that the members are inside the resulting group. EVoC does the same
thing more politely ($0.557$ recall, $0.003$ precision). This is the trap that
the literature's habit of quoting a single number walks into, and it is why
\textbf{Table~\ref{tab:field}} prints three.

\textbf{Spectra read directly beat element ratios, and it took a chemical
ablation to be sure.} The 64-dimensional PCA latent and the 256-dimensional
masked-AE latent both reach recall $\approx 0.42$ at precision
$\approx 0.62$ from the abundances' $0.20$/$0.24$. Two independent
representations of the same spectra agree with each other to within
$0.002$ in recall and $0.012$ in precision --- which is the honest way to say
that the improvement is real while the difference between the $64$- and
$256$-dimensional spectral rows is not.

\textbf{The gap to the ceiling is the story.} Kinematics alone recover $0.93$
of the members at a precision limited by velocity doppelg\"angers rather than
by the method (\S\ref{sec:fundamentals}). Chemistry, read as well as we can
read it today, gets to $0.42$ recall. Strong chemical tagging --- recovering a
dispersed cluster with no kinematic help --- is not what these numbers
demonstrate; what they demonstrate is that modern self-supervised spectral
representations find \emph{something} real about a cluster in a spectrum, and
that the something is diluted by a factor of two relative to parallaxes and
proper motions.

\subsection{The curse of dimensionality, measured}

Concatenating a high-dimensional, weakly informative signal onto a
low-dimensional, strongly informative one does not add their information. It
dilutes the strong signal, because distance is shared across dimensions. The
DR17 study measured this on 878 members in region mode with the kinematic
labels held out of the ground truth:

\begin{table}[htbp]
\centering
\small
\tabcolsep5pt
\begin{tabular}{@{}lccc@{}}
\toprule
signal (dims) & cluster homogeneity & field recall & field precision \\
\midrule
abundances (16)          & 0.25 / 0.51 / 0.51 & 0.73 & 0.47 \\
spectral latent (256)    & 0.40 / 0.69 / 0.64 & 0.72 & 0.51 \\
+ kinematics (20 / 260)  & 0.62 / 0.74 / 0.72 & 0.83 / 0.80 & 0.63 / 0.61 \\
kinematics only (4)      & \textbf{0.97 / 0.96 / 0.92} & \textbf{0.93} & 0.63--0.67 \\
\bottomrule
\end{tabular}
\caption{Signal comparison on the DR17 data (three numbers per cell are the
t-SNE / UMAP / EVoC homogeneity). Kinematics alone near-solve the cluster-only
task and set the field-recall ceiling; adding chemistry on top improves the
cluster-only score slightly and never improves recall beyond what kinematics
already achieved. The DR19 re-run reproduces the pattern with different
absolute values (kinematics $0.94$, abundances $0.56$,
abundances$\,+\,$kinematics $0.77$ over the same 982 stars).}
\label{tab:curse}
\end{table}

The mechanism is worth stating precisely, because it is the reason this
workbook spends three chapters on geometry rather than one on algorithms. A
distance in a 20-dimensional space is an average over 20 coordinate
differences. If 16 of those coordinates are chemistry and carry weak
information about cluster identity, while 4 are kinematics and carry strong
information, then the strong coordinates are outvoted$\,$--- not mis-weighted,
outvoted, because no reweighting recovers information that the metric has
already blurred. Rescaling chemistry \emph{up} to compensate costs precision
in the field, as \S\ref{sec:benchmark}'s lever table shows. The correct
response is architectural, not numerical: cluster on the strong signal, then
apply the weak signal as a rejection test. That is the two-stage pipeline, and
it is the single most successful idea in this project.

\subsection{Levers, and a lesson about bug fixes}

The DR17 study swept the pipeline's knobs on one cluster (M~67) inside a
$30^\circ$ region, with $5\,000$ stars and 183 referee members. Every run
logged its full configuration. The table is worth reading as a list of
\emph{hypotheses}, not results, and the section that follows explains why.

\begin{table}[htbp]
\centering
\small
\tabcolsep4pt
\begin{tabular}{@{}lcccc@{}}
\toprule
lever & recall & prec. & UMAP recall & verdict \\
\midrule
baseline                                   & 0.087 & 0.117 & 0.415 & --- \\
\texttt{min\_cluster\_size} $5\!\to\!10$   & 0.404 & 0.129 & 0.525 & best single \\
perplexity $30\!\to\!50$                   & 0.383 & 0.124 & 0.415 & free win \\
perplexity $30\!\to\!10$                   & 0.033 & 0.353 & 0.415 & precision for recall \\
\texttt{SNR\_MIN} $100\!\to\!150$          & 0.296 & 0.089 & 0.310 & recall up, precision down \\
dwarf-only sample                          & 0.149 & 0.165 & 0.158 & best purity \\
drop S, K, V, Mn                           & 0.055 & 0.256 & 0.038 & hurts recall \\
UMAP \texttt{metric=cosine}                & 0.087 & 0.117 & 0.044 & hurts UMAP \\
UMAP \texttt{n\_neighbors} $=5$            & 0.087 & 0.117 & 0.071 & worse \\
UMAP \texttt{n\_neighbors} $=30$           & 0.087 & 0.117 & 0.410 & $=$ baseline \\
element weights (after a bug fix)          & \textbf{0.437} & \textbf{0.158} & 0.339 & best lever \\
weights $+$ size floor                     & 0.432 & 0.160 & 0.525 & not additive \\
\bottomrule
\end{tabular}
\caption{DR17 lever sweep (M~67, $30^\circ$ region, $5\,000$ stars, 183
referee members, seed 42). The recall and precision columns are t-SNE; the
last two rows exist because of a bug: the element-weighting switch was
originally a no-op, since the weights were applied and then undone by the
standardisation step that followed. Reordering the two operations turned the
largest apparent non-effect into the largest effect in the table. Two rules of
thumb survive from this sweep --- the size floor of $5$ fragments the field,
and a perplexity below $N/100$ starves the embedding --- and one warning: the
best levers do not compose.}
\label{tab:levers}
\end{table}

The lesson is not that hyperparameters matter, which is known, but that the
\emph{largest single effect in a sweep can be a bug}. A configuration flag that
does nothing is worse than a missing flag, because it produces a negative
result that looks like an experiment. The DR17 re-check of the same lever on
DR19 data did not reproduce the recall gain, which is exactly what the protocol
of \S\ref{sec:validation} predicts for a lever that was selected on one
cluster at one seed: selecting on a metric and reporting the same metric is
not validation. The DR19 region sweep of the same switch confirms the warning
from the other direction --- a recall that climbs to $0.96$ while precision
falls to $0.05$ is a cluster that has grown to swallow the field.

\subsection{Where this sits in the literature: Casamiquela et al. (2021)}
\label{sec:literature}

Every number so far in this chapter has been scored against our own baselines.
The closest published measurement is \citet{Casamiquela:21}, and it is close
enough to sit in the same table: the same metric triple, the same task, and a
sample built to be the best case that spectroscopy can produce. They measured 16
elements differentially in 175 red-clump members of 31 thin-disc open clusters,
with an internal coherence of about $0.03$ dex per element and \emph{no field
stars}. They then clustered the abundances directly with HDBSCAN --- no
embedding, no projection --- and fine-tuned its free parameters to maximise the
number of correctly recovered clusters (their Sect. 4.2). They call the sample
itself their \emph{best-case scenario}. Their published result is $h = 0.49$, $c = 0.63$, $V = 0.55$,
$\mathrm{RF}_{40} = 0.29$ (9 of 31 clusters) and $\mathrm{RF}_{70} = 0.03$ (one
cluster, NGC~2682); their own summary of that run is that more than 70\,\% of the
groups it finds are statistical, mixing stars from several real clusters.

\begin{table}[htbp]
\centering
\small
\tabcolsep5pt
\begin{tabular}{@{}lcccccc@{}}
\toprule
sample & $n_\star$ / $n_{\rm cl}$ & $h$ & $c$ & $V$ & $\mathrm{RF}_{40}$ & $\mathrm{RF}_{70}$ \\
\midrule
\citet{Casamiquela:21}, published                  & 175 / 31 & 0.49  & 0.63  & 0.55  & 0.290 & 0.030 \\
their clustering step, our stars (raw dex)         & 982 / 25 & 0.499 & 0.416 & 0.454 & 0.080 & 0.040 \\
their clustering step, our stars (standardised)    & 982 / 25 & 0.475 & 0.407 & 0.438 & 0.000 & 0.000 \\
\midrule
our pipeline: t-SNE                                & 982 / 25 & 0.560 & 0.601 & 0.580 & 0.320 & 0.080 \\
our pipeline: UMAP                                 & 982 / 25 & 0.578 & 0.572 & 0.575 & 0.326 & 0.006 \\
our pipeline: EVoC                                 & 982 / 25 & 0.420 & 0.529 & 0.466 & 0.189 & 0.006 \\
our pipeline: K-means ($K$ known)                  & 982 / 25 & 0.454 & 0.490 & 0.471 & 0.080 & 0.000 \\
\midrule
open clusters only, UMAP                           & 665 / 18 & 0.573 & 0.422 & 0.486 & 0.405 & 0.040 \\
open clusters only, t-SNE                          & 665 / 18 & 0.471 & 0.458 & 0.465 & 0.222 & 0.000 \\
open clusters only, EVoC                           & 665 / 18 & 0.520 & 0.419 & 0.464 & 0.214 & 0.048 \\
\bottomrule
\end{tabular}
\caption{The published comparison. Rows 2--3 are \textit{their} clustering step
--- HDBSCAN on the abundances, \texttt{min\_cluster\_size}$=2$, their
$4 \times 7$ parameter grid, best grid point by $\mathrm{RF}_{40}$ --- run on
our stars; their configuration is defined in dex, so the raw-abundance row is
the faithful replication and the standardised row is a rescaling check. Their
$4 \times 7$ grid never reaches $\mathrm{RF}_{40} > 0.11$ on our stars. Rows
4--10 are our arms, mean over seven seeds (the $h$, $c$, $V$ columns repeat
\textbf{Table~\ref{tab:honest}}); the open-cluster subset is the closest
analogue of their sample. Chance levels, from shuffling the true labels into
groups of the observed sizes: for their sample shape $V = 0.51$ ($0.44$--$0.51$
across the three partition shapes tested) and $\mathrm{RF}_{40} = 0.04$ (95th
percentile $0.10$); for their clustering step run on \emph{our} stars $V = 0.29$
--- their rule cuts 165 groups, and many small groups inflate the null --- with
$\mathrm{RF}_{40} \le 0.001$; for our own arms $V = 0.06$--$0.13$ and
$\mathrm{RF}_{40} \le 0.001$. Source:
\texttt{results/casamiquela\_comparison*.csv} and
\texttt{results/casamiquela\_protocol.csv}; the second one also carries the grid
search.}
\label{tab:literature}
\end{table}

Four things follow from that table.

\textbf{Their clustering step recovers less on our stars than on theirs.} Run
their configuration on our 982 stars and the recovery fraction falls to
$0.08$ ($0.11$ on the open clusters alone, both being Berkeley~17 and
NGC~1245), against $0.29$ on their own sample; under the standardised variant it
returns 168 groups and recovers nothing at all. Two differences remain that we
cannot separate: our abundances are survey-quality --- $\sim\!0.05$ dex per
element against their $0.03$ --- and our clusters are not a single evolutionary
stage.

\textbf{Our two-stage pipeline lands in the same range as their one-stage
pipeline.} t-SNE and UMAP on 16 survey abundances, with default HDBSCAN, recover
$0.32$ and $0.33$ of the 25 clusters, and on the 18 open clusters alone --- the
subset closest to their sample --- UMAP reaches $0.41$. The honest reading is not
that we beat them. It is that embed-then-cluster on coarser abundances reaches
the same recovery as cluster-the-abundances on the best abundances available,
and that neither gets past about a third of the clusters.

\textbf{The failures have the same shape.} Their surviving clusters sit at the
edges of the abundance distribution --- NGC~6705, NGC~2682, and NGC~2420, which
they call the most metal-poor cluster in their sample. NGC~2420 is a regular
recovery of ours too: UMAP recovers it in every seed of the open-cluster subset
and EVoC in six of seven. Our arms otherwise recover the same handful of
clusters in every seed --- Berkeley~17, IC~166, M~71, NGC~1245, NGC~6819,
NGC~7789, the Pleiades --- and that set contains M~3, which \S\ref{sec:clusters}
introduces as the metal-poor outlier of the sample. Whether this is the same
edge-of-distribution mechanism their paper describes is a hypothesis this
workbook does not settle: member evolutionary state is a confound we have not
yet removed, which is what the \texttt{DWARF\_ONLY} lever of
\S\ref{sec:assignment} is for.

\textbf{Their field-star test is the published analogue of our Task 2, and it
agrees.} Adding field stars to their sample leaves one cluster above the 40\,\%
threshold (NGC~188) and none at 70\,\%. Our field retrieval
(\textbf{Table~\ref{tab:field}}) agrees in direction and goes further: the best
abundance-only arm reaches recall $0.20$ at precision $0.24$, and the arms that
reach high recall do it by labelling most of the field as cluster. Their element
tests agree with ours too: restricting the set to 8 or 7 elements did not improve
their recovery, and our dropped-element row in \textbf{Table~\ref{tab:levers}}
loses recall as well.

\subsection{What this benchmark cannot tell you}
\label{sec:limitations}

Five limits, in descending order of how much they should trouble a reader.

\begin{enumerate}
\item \textbf{The sample is small and uneven.} 982 stars across 25 clusters is
  a few tens of stars per object, and the distribution is wide: M~67 supplies
  218 stars and the Pleiades 113, while five clusters supply fewer than ten each
  once the duplicate rows are collapsed (Berkeley~17, NGC~1798, King~5,
  NGC~2158, King~7 --- \textbf{Table~\ref{tab:clusters}} counts rows, and lists
  three). \textbf{Table~\ref{tab:ablation}}
  shows what a single well-populated cluster can do to a result, and every macro
  average in this chapter is therefore an average over a handful of clusters
  that are not equally measured. The member frame is itself
  configuration-dependent: the 1\,002 rows counted in
  \textbf{Table~\ref{tab:clusters}} need a finite-abundance floor of nine out
  of sixteen elements, while the shipped default of eight keeps 1\,215 rows ---
  so the saved baseline scores are not reproducible from a bare
  \texttt{cluster baseline} run
  (\texttt{python scripts/population\_counts.py --baseline-min-finite 9}).
\item \textbf{The member matrix contains duplicated stars.} 298 rows share an
  \texttt{APOGEE\_ID} with another row, the duplicates differ by
  $\sim\!0.04$ dex per element, and they inflate the largest clusters most.
  Until the sample is collapsed to unique stars, every macro average in this
  chapter is a weighted average over something that is not quite a
  star.
\item \textbf{Some rows are not measurements.} The abundances-only t-SNE row
  is row-order unstable (\S\ref{sec:validation}); the ablation row is
  degenerate by construction. Both are printed rather than dropped, because
  the pattern of which arms are fragile is itself a result: the
  spectral-latent arms are stable, and the raw-abundance t-SNE arm is not.
\item \textbf{The external comparison is one pipeline, not a survey of the
  literature.} \S\ref{sec:literature} runs \citet{Casamiquela:21}'s clustering
  step on our stars and their metric triple on our scores, which is as far as a
  like-for-like comparison goes without their star list. We do not re-run the
  graph-attention autoencoder of \citet{Spina:25}, which works on
  $\sim\!47\,000$ APOGEE stars and deliberately adds kinematics and age to the
  chemistry, so any claim of state-of-the-art in this space should be read as
  unaudited.
\item \textbf{The counter-result stands, and it is not the same measurement.}
  \citet{Hogg:16} recovered phase-space structure by clustering 15 abundances
  with K-means, and our K-means row in \textbf{Table~\ref{tab:honest}} matches
  EVoC's homogeneity. On the recovery fraction it is the weakest arm of the
  table ($0.08$ against t-SNE's $0.32$), which is the difference between a tidy
  partition and a recovered cluster: $K$-means is told how many groups to make,
  so its homogeneity is high by construction, while the density methods have to
  find the groups. Whatever this workbook demonstrates about the limits of
  abundance-space clustering, it is a statement about APOGEE's ASPCAP
  precision, these 25 clusters, and these methods --- not a statement that
  chemical tagging does not work.
\end{enumerate}

\begin{issues}[EXERCISES]
\begin{enumerate}
\item Reproduce one row of \textbf{Table~\ref{tab:honest}}: pick a signal and
  a method, run it over the seven seeds, and report mean $\pm$ standard
  deviation. Compare with the published row. If your numbers differ, is the
  cause the population, the seed set, or a default?
\item \textbf{Table~\ref{tab:field}} reports precision for the kinematic
  referee. Recompute it against the Simbad memberships and plot the
  difference per cluster. Which clusters are most referee-sensitive, and does
  that correlate with their age or their [Fe/H]?
\item The ablation of \textbf{Table~\ref{tab:ablation}} removes M~3. Repeat
  it for the other clusters with more than 100 rows (M~67, NGC~6819, NGC~2243)
  and for the duplicates-cleaned matrix. How much of the instability survives
  cleaning?
\item \S\ref{sec:literature} runs their pipeline on our stars but not ours on
  theirs. Their sample reachable star list is not public in a form we can
  re-run, so state what you would need from the other authors --- and what a
  re-run on their stars would change about the conclusion of
  \textbf{Table~\ref{tab:literature}}.
\end{enumerate}
\end{issues}
